%! Author = chtompki
%! Date = 1/14/26
\documentclass[12pt]{amsart}
% Default margins are too wide all the way around. I reset them here
\setlength{\hoffset}{-.75in}
\setlength{\textwidth}{6.5in}
\setlength{\voffset}{-.5in}
\setlength{\textheight}{9.0in}
\usepackage{amsmath}
\usepackage{amssymb}
\usepackage{amsthm}
\usepackage{enumitem}
\usepackage{setspace}
\usepackage{graphicx}
\usepackage{hyperref}

\renewcommand{\thesubsection}{\thesection.\alph{subsection}}

\newenvironment{solution}
{\begin{proof}[Solution]}
{\end{proof}}

\NewDocumentCommand{\vitem}{v}{%
    \item[\texttt{#1}]%
}

\title{Numerical Analysis\\Problem Set 3}
\author{Rob Tompkins}
\date{\today}

\begin{document}
    \maketitle
    \date{\today}
    \begin{onehalfspace}
        We begin by noting that all of the files referenced in this document are available at: \\
        \underline{https://github.com/chtompki/MathSpring2026/tree/main/NumericalAnalysis/ProblemSet3}

        \section{Problem 1}
        We'll be approximating $\sqrt{2}$ though this problem

        %1.a
        \subsection{}
        Our goal is to find a polynomial interpolant $p(x)$ using the Lagrange form given that we have
        three points in the form $x,f(x)$ as $\{(-1,1/2), (0,1), (1,2)\}$. So we let $f(x) = 2^x$ such that
        $f(1/2) = \sqrt{2}$. So our polynomial begins as
        \[
            p(x) = y_0 L_0(x) + y_1 L_1(x) + y_2 L_2(x),
        \]
        with $y_i = f(x_i)$ and $L_i(x)$ being the $i^{th}$ Lagrange polynomial.
        Computing the polynomials follows as
        \begin{gather*}
            L_{2,0}(x) = \frac{(x-x_1)(x-x_2)}{(x_0-x_1)(x_0-x_2)} = \frac{(x-0)(x-1)}{(-1-0)(-1-1)} = \frac{x(x-1)}{2}\\
            L_{2,1}(x) = \frac{(x-x_0)(x-x_2)}{(x_1-x_0)(x_1-x_2)} = \frac{(x+1)(x-1)}{(0+1)(0-1)} = 1-x^2\\
            L_{2,2}(x) = \frac{(x-x_0)(x-x_1)}{(x_2-x_0)(x_2-x_1)} = \frac{(x+1)(x-0)}{(1+1)(1-0)} = \frac{x(x+1)}{2}
        \end{gather*}

        %1.b
        \subsection{}
        Putting these together we have the following polynomial:
        \begin{eqnarray*}
            p(x) & = &  \frac{1}{2}\cdot\frac{x(x-1)}{2} + 1\cdot(1-x^2) + 2\cdot\frac{x(x+1)}{2} \\
             & = & \frac{x^2}{4} - \frac{x}{4} + 1 -x^2 + x^2 + x\\
            & = & \frac{x^2}{4} + \frac{3x}{4} + 1
        \end{eqnarray*}
        Pluggin in $x=1/2$ we get the following
        \begin{eqnarray*}
            p\left(\frac{1}{2}\right) & = & \frac{1}{4}\cdot\frac{1}{4} + \frac{3}{4}\cdot\frac{1}{2} + 1\\
            & = & \frac{1}{16} + \frac{3}{8} + 1\\
            & = & \frac{1}{16} + \frac{6}{16} + \frac{16}{16}\\
            & = & \frac{23}{16} \approx \sqrt{2}.
        \end{eqnarray*}

        %1.c
        \subsection{}
        For the error bound we want to compute the following
        \begin{displaymath}
            |f(x) - p(x)| = \frac{f'''(\xi)}{3!} (x - x_0)(x - x_1)(x - x_2),
        \end{displaymath}
        where $\xi\in[-1,1]$.
        We compute the third derivative of $f(x)$ by setting $f(x) = \ln(x) = e^{x\ln(2)}$.
        Clearly then we have
        \begin{displaymath}
            f'''(x) = (\ln(2))^3 e^{x\ln(2)} = (\ln(2))^3 2^x.
        \end{displaymath}
        We want to have $\xi$ such that $|f'''(\xi)|$ is largest so that everything else is beneath that
        as our upper bound.
        Because $f(x) = 2^x$ is an increasing exponential, the maximum of the third derivative will necessarily
        be the right hand end point. So we set $\xi=1$ giving us $f'''(\xi) = 2(\ln(2))^3.$ So
        \begin{displaymath}
            |f(x) - p(x)| = \frac{f'''(\xi)}{3!} (x - x_0)(x - x_1)(x - x_2) = \frac{2(\ln(2))^3}{6}(x+1)(x)(x-1)
        \end{displaymath}
        And we want to know the error at $x=1/2$ which comes out at
        \begin{eqnarray*}
            \left|f\left(\frac{1}{2}\right) - p\left(\frac{1}{2}\right)\right| & = & \frac{2(\ln(2))^3}{6}\left(\frac{1}{2}+1\right)\cdot\frac{1}{2}\cdot\left(\frac{1}{2}-1\right)\\
            & = & \frac{3\ln(2)^3}{24} \\
            & = & \frac{\ln(2)^3}{8} = 0.0416
        \end{eqnarray*}

        \section{Problem 2}
        Our goal here is to approximate $\ln(2)$ with piecewise polynomials.

        %2.a
        \subsection{}
        Since each piece has 3 coefficients and there are $n+1$ points there are $n$ pieces to the polynomial.
        So we have $3n$ coefficients to account for.
        We will refer to our interpolation function at $s(x)$, with the piceces being $s_i(x)$.

        We begin by interpolating the endpoints of the segments giving us $s_i(x_i)=y_i$ and $s_i(x_{i+1}) = y_{i+1}$.
        This gives us our first $2n$ constraints.

        Next we approach the first derivatives of the pieces for which we need equality for continuity on the
        inner $n-1$ points or knots.
        So we must have that
        \[
            s_i'(x_{i+1}) = s_{i+1}(x_{i+1}),
        \]
        which gives us another $n-1$ conditions.
        So we are missing 1 condition to reach our $3n$ requisite conditions to be able to solve for

        The piece that we are missing is one of the boundary conditions. For example we could use
        $s_0(x_0)$ and fix that or use $s_{n-1}(x_n)$.
        Regardless we need this one more condition to get all of our constants.

        %2.b
        \subsection{}
        We were instructed to use the dataset from problem 1, $\{(-1,1/2), (0,1), (1,2)\}$ to solve this problem.
        This yields the intervals $I_0=[-1,0]$ and $I_1=[0,1]$.
        The interpolating polynomials that we are looking to find the coefficients for follow
        \begin{gather*}
            s_0(x) = a_0 + b_0(x+1)+c_0(x+1)^2\\
            s_1(x) = a_1 + b_1 x + c_1 x^2.
        \end{gather*}
        We begin by using the interpolation conditions such that $s_0(-1) = 1/2 \Rightarrow a_0=1/2$,
        $s_0(0) = 1 = 1/2 + b_0 + c_0$ and $s_0'(-1) = b_0 + 2c_0(-1+1) = 0$ we have that $b_0 = 0$.
        Thus $c_0 = 1/2$, which gives us
        \[
            s_0(x) = \frac{1}{2} + \frac{1}{2}(x+1)^2
        \]
        Notice that $s_1(0) = 1 \Rightarrow a_1 = 1$.
        Next we use derivative continuity at $x=0$ to find the other coefficients.
        We have that $s_1'(x) = b_1 + 2c_1 x$, so $s_0'(0) = s_1'(0) = b_1 = 1$.
        Lastly, we use our boundary condition $s_1(1) = 2 = 1 + 1 + c_1(1) = 2$ to get $c_1 = 0$.
        And we have
        \[
            s_1(x) = 1 + x.
        \]
        So
        \[
            s(x) =
            \begin{cases}
                \frac{1}{2} + \frac{1}{2}(x+1)^2 & -1\le x < 0\\
                1+x & 0\le x \le 1
            \end{cases}
        \]

        %2.c
        \subsection{}
        We begin by noticing that $s_0'(0) = 1$ and $s_1'(0) = 1$ by the continuity of $s'(x)$.
        Now again we are using the function $f(x)=2^x$ such that $f'(x) = \ln(2) 2^x$ and $f'(0) = \ln(2)$, so
        we use $s'(0) = 1 \approx \ln(2)$.
        We did a calculation on a computer just to test the error here and we got $\ln(2)=0.69315$, so we
        are quite off

        %2.d
        \subsection{}
        Skip

        \section{Problem 3}


    \end{onehalfspace}
\end{document}